\documentclass[aspectratio=169]{beamer}
\input{../../lib/preamble}

\title{ML Foundations}
\subtitle{\coursetitle}
\author{\instructor}
\date{\today}

\begin{document}

\begin{frame}
\titlepage
\end{frame}

\begin{frame}{Outline}
\tableofcontents
\end{frame}

% ============================================================
\section{The Training Loop}
% ============================================================

\begin{frame}{Concepts in Machine Learning}
Four ingredients:
\begin{description}
\item[Data] some input to a task we want an AI to do
\item[Model] gets data, outputs something (a classification, an action, a prediction)
\item[Loss function] quantifies how well the model did
\item[Optimizer] updates the model to do the task better
\end{description}
\end{frame}

\begin{frame}[fragile]{The Training Loop in PyTorch}
\begin{minted}{python}
for x, y in dataloader:
    optimizer.zero_grad()
    prediction = model(x)         # forward pass
    loss = loss_fn(prediction, y) # compute loss
    loss.backward()               # backward pass
    optimizer.step()              # update parameters
\end{minted}
\vspace{1em}
This is \emph{the} loop.
Everything else in ML is about making each piece better.
\end{frame}

\begin{frame}{Terminology}
\begin{description}
\item[Parameters / weights] the numbers that represent the model itself, updated by the optimizer ($\theta$)
\item[Activations] the numbers created as we feed data through the network
\item[Hyperparameters] choices made \emph{before} training: learning rate, batch size, architecture
\end{description}
\end{frame}

\begin{frame}{Stochasticity and Batching}
We want to minimize the true loss over the full data distribution:
$$\mathcal{L}(\theta) = \mathbb{E}_{x \sim \mathcal{D}} \left[ \ell(f_\theta(x), y) \right]$$
But we only ever see a finite \textbf{batch} $\mathcal{B} \subset \mathcal{D}$ at each step:
$$\hat{\mathcal{L}}(\theta) = \frac{1}{|\mathcal{B}|} \sum_{(x,y) \in \mathcal{B}} \ell(f_\theta(x), y)$$
\begin{itemize}
\item Every gradient is a noisy estimate of the true gradient
\item Batch size controls the noise level
\end{itemize}
\end{frame}

% ============================================================
\section{Tensors and PyTorch}
% ============================================================

\begin{frame}[fragile]{Tensors}
The central object in PyTorch: multidimensional arrays of numbers.
\vspace{0.5em}

We use tensors to represent: data, parameters, activations, outputs, losses.
\vspace{0.5em}

\begin{minted}{python}
x = torch.tensor([1.0, 2.0, 3.0])
M = torch.zeros(3, 4)
\end{minted}
\vspace{0.5em}

By convention, everything except parameters has the \textbf{batch dimension first}, so we can process many samples at once.
\end{frame}

\begin{frame}[fragile]{Tensor Operations}
\begin{minted}{python}
torch.exp(x)           # elementwise exponential
A @ B                  # matrix multiply
x.unsqueeze(0)         # add a dimension
\end{minted}
\vspace{1em}

Einstein notation makes index operations explicit:
$$C_{ij} = \sum_k A_{ik} B_{kj}$$
\begin{minted}{python}
torch.einsum('ik,kj->ij', A, B)
\end{minted}
\end{frame}

\begin{frame}[fragile]{Einops}
Readable notation for tensor reshaping:
\begin{minted}{python}
# PyTorch (what does this do?)
x.reshape(B, S, n_heads, d_head).permute(0, 2, 1, 3)

# Einops (self-documenting)
rearrange(x, 'b s (h d) -> b h s d', h=n_heads)
\end{minted}
\vspace{0.5em}
\begin{itemize}
\item \texttt{rearrange}: reshape, permute, split, merge
\item \texttt{reduce}: pool over dimensions
\item \texttt{repeat}: tile/broadcast
\end{itemize}
\end{frame}

\begin{frame}{Autograd and the Computation Graph}
Tensors keep track of a computational graph as you operate on them.
\vspace{0.5em}
\begin{enumerate}
\item Build up expressions from tensors with \texttt{requires\_grad=True}
\item Call \texttt{.backward()} on a scalar result
\item Every tensor in the graph can now access its gradient via \texttt{.grad}
\end{enumerate}
\vspace{1em}
You never write derivatives by hand.
The graph is rebuilt every forward pass (``define-by-run'').
\end{frame}

\begin{frame}[fragile]{Devices}
\begin{itemize}
\item Tensors live on a \textbf{device}: CPU or GPU (\texttt{cuda:0})
\item Move with \texttt{x.to('cuda')}
\item All tensors in an operation must be on the same device
\item GPU memory is the main bottleneck for model size and batch size
\end{itemize}
\end{frame}

\begin{frame}{}
\centering
\Large \textbf{Exercise 1: PyTorch Basics} \\[1em]
\normalsize 25 minutes
% TODO: add Colab link and QR code
\end{frame}

% ============================================================
\section{Loss Functions}
% ============================================================

\begin{frame}{Loss Functions: Supervised Learning}
The loss function defines \emph{what} we are optimizing.
\vspace{0.5em}

When we can directly verify the output:
\begin{itemize}
\item Predict something continuous (position of a ball) $\Rightarrow$ MSE
\item Predict something discrete (class, next word) $\Rightarrow$ \textbf{cross-entropy loss}
$$\mathcal{L} = -\sum_i y_i \log \hat{y}_i$$
\end{itemize}
\vspace{0.5em}
Minimizing cross-entropy makes the model \textbf{well-calibrated}: its output numbers can be interpreted as probabilities.
\end{frame}

\begin{frame}{Loss Functions: Human Preferences}
We do not always have the ``correct label.''
Sometimes we only know a pattern outputs should have.
\vspace{0.5em}

Example: rating text quality, given only pairwise human preferences (``I prefer A over B'').
\vspace{0.5em}

The preferred text should be rated higher:
$$\mathcal{L} = -\log \sigma(r_\theta(x_w) - r_\theta(x_l))$$
\cite{christiano2023a}
\end{frame}

\begin{frame}{Loss Functions: Reinforcement Learning}
Sometimes we cannot say what the correct output looks like, but we can identify one when we see it (e.g.\ winning at chess).
\vspace{0.5em}

\begin{itemize}
\item RL loss functions are more complex
\item Often involve a learned reward model (from the previous slide) combined with a policy objective
\end{itemize}
\end{frame}

\begin{frame}{Train Loss vs.\ Test Loss}
% TODO: add overfitting/underfitting plot
\vspace{1em}
\begin{description}
\item[Underfitting] high train loss, high test loss (model too simple or undertrained)
\item[Overfitting] low train loss, high test loss (model memorizes training data)
\item[Good fit] both low (generalizes to unseen data)
\end{description}
\vspace{1em}
Regularization (L1, L2, dropout) constrains the model to reduce overfitting.
\end{frame}

% ============================================================
\section{Optimizers}
% ============================================================

\begin{frame}{Optimizers}
We have a gradient.
Which direction do we step?
\end{frame}

\begin{frame}{Stochastic Gradient Descent}
$$\theta_{t+1} = \theta_t - \alpha \, \nabla_\theta \mathcal{L}$$
\begin{itemize}
\item $\alpha$: learning rate (hyperparameter)
\item Simple and robust
\item Can be slow: oscillates across narrow valleys, crawls along them
\end{itemize}
\vspace{1em}
% TODO: add picture of descent with too-small and too-large learning rate
\end{frame}

\begin{frame}{Momentum}
$$v_{t+1} = \beta \, v_t + (1 - \beta) \, \nabla_\theta \mathcal{L}$$
$$\theta_{t+1} = \theta_t - \alpha \, v_{t+1}$$
\begin{itemize}
\item Keeps a running velocity, like a ball rolling downhill
\item Accelerates through flat regions
\item Dampens oscillations in steep directions
\item Can roll over small local minima
\end{itemize}
\vspace{1em}
% TODO: add picture of momentum rolling over a local minimum
\end{frame}

\begin{frame}{Adam = Momentum + RMSProp}
$$m_{t+1} = \beta_1 m_t + (1-\beta_1) g_t \qquad v_{t+1} = \beta_2 v_t + (1-\beta_2) g_t^2$$
$$\hat{m} = \frac{m_{t+1}}{1-\beta_1^{t+1}} \qquad \hat{v} = \frac{v_{t+1}}{1-\beta_2^{t+1}}$$
$$\theta_{t+1} = \theta_t - \alpha \frac{\hat{m}}{\sqrt{\hat{v}} + \epsilon}$$
\begin{itemize}
\item Adaptive per-parameter learning rates + momentum
\item Default choice for most deep learning
\item Typical: $\beta_1=0.9$, $\beta_2=0.999$, $\epsilon=10^{-8}$
\end{itemize}
\end{frame}

\begin{frame}{}
\centering
\Large \textbf{Exercise 2: Implement Optimizers} \\[1em]
\normalsize 25 minutes
% TODO: add Colab link and QR code
\end{frame}

% ============================================================
\section{Architectures}
% ============================================================

\begin{frame}{The MLP}
% TODO: add diagram of a simple MLP
\vspace{1em}
\begin{itemize}
\item Input $\rightarrow$ hidden layer $\rightarrow$ output
\item Without a hidden layer, we can only represent linear functions
\item The hidden layer with a nonlinearity lets us express richer functions
\end{itemize}
\end{frame}

\begin{frame}{Activation Functions}
Without nonlinearities, stacking linear layers just gives another linear layer.
\vspace{0.5em}
\begin{itemize}
\item \textbf{ReLU}: $\max(0, x)$, simple and effective
\item \textbf{GELU}: smooth approximation, used in transformers
\item \textbf{Sigmoid/Tanh}: saturate at extremes, mostly historical
\end{itemize}
\end{frame}

\begin{frame}{Universal Approximation}
\textbf{Theorem} (Cybenko, Hornik et al.):
A single hidden layer with nonlinear activations can approximate any continuous function on a compact set to arbitrary precision.
\vspace{1em}
\begin{itemize}
\item Existence result, not a practical recipe
\item Says nothing about how many neurons you need
\item In practice, deeper is much more efficient
\end{itemize}
\end{frame}

\begin{frame}{Depth: Stacking Layers}
We can more easily compute things if we compute them in multiple steps.
\vspace{0.5em}
\begin{itemize}
\item Example: classifying a dog is easier if early layers detect ears and snouts, later layers combine them into ``head''
\item Models have layers that mostly look the same but learn different things
\item The deep learning recipe: repeat a simple block many times
\end{itemize}
\vspace{0.5em}
Over/underparameterization:
\begin{itemize}
\item Too few parameters $\Rightarrow$ cannot represent the function
\item Too many $\Rightarrow$ can memorize, but modern networks often still generalize
\end{itemize}
\end{frame}

\begin{frame}{Residual Connections}
Deep networks suffer from \textbf{vanishing gradients}: the chain rule often makes gradients go to zero through many layers.
\vspace{0.5em}

Fix: add a skip connection so every layer can start learning something.
$$x_{l+1} = x_l + f(x_l)$$
\begin{itemize}
\item Gradient through the skip is exactly 1
\item Network only needs to learn the \emph{residual} $f(x_l)$
\item Enables training networks with hundreds of layers
\end{itemize}
\vspace{0.5em}
% TODO: add diagram of residual stream with blocks feeding in/out
\end{frame}

\begin{frame}{Architecture Encodes Symmetry}
The right architecture exploits structure in the data.
\vspace{0.5em}

Example: in an image, a dog ear looks the same regardless of position.
We need a detector that is \textbf{translation invariant}, so we should not learn a separate detector for every location.
\vspace{1em}
An architecture that respects a symmetry needs fewer parameters and less data to learn the same function.
\end{frame}

\begin{frame}{Convolutional Neural Networks}
\begin{itemize}
\item A small kernel slides across the input (weight sharing)
\item Same features detected regardless of position
\item Hierarchy: edges $\rightarrow$ textures $\rightarrow$ parts $\rightarrow$ objects
\end{itemize}
\vspace{1em}
% TODO: add gif/animation of convolution sliding over an image
\end{frame}

\begin{frame}{Transformers (Preview)}
Text has a similar symmetry: each word is similar, but understanding requires long-range connections, not just neighbors.
\vspace{0.5em}
\begin{itemize}
\item \textbf{Self-attention}: each token attends to all other tokens
\item Repeated blocks of: attention $\rightarrow$ feedforward $\rightarrow$ residual connection
\item Dominant architecture for language (and increasingly vision, audio, \ldots)
\item We will go much deeper in the afternoon session
\end{itemize}
\vspace{0.5em}
% TODO: add gif/animation of transformer attention
\end{frame}

\begin{frame}{}
\centering
\Large \textbf{Exercise 3: Build Simple Architectures} \\[1em]
\normalsize 15 minutes
% TODO: add Colab link and QR code
\end{frame}

% ============================================================
\section{Hyperparameter Optimization}
% ============================================================

\begin{frame}{Hyperparameters}
Any property not set by gradient descent is set by the experimenter:
\begin{itemize}
\item Model architecture (depth, width, \ldots)
\item Batch size, number of epochs
\item Learning rate, momentum, and other optimizer settings
\end{itemize}
\end{frame}

\begin{frame}{Finding Good Hyperparameters}
\begin{itemize}
\item Some we have figured out once and reuse (e.g.\ $\beta_1 = 0.9$)
\item Some we sweep over: run many times, pick the best (typically learning rate)
\item Some follow \textbf{scaling laws}: cheap small experiments predict large-scale performance
\end{itemize}
\end{frame}

\begin{frame}{Scaling Laws}
$$L(N) \propto N^{-\alpha}$$
\begin{itemize}
\item Loss decreases as a power law with model size, data, or compute
\item Enables predicting performance before training the full model
\item Compute-optimal training: balance model size and data % TODO: add Chinchilla citation
\end{itemize}
\vspace{1em}
% TODO: add Chinchilla scaling law plot
\end{frame}

\begin{frame}{}
\centering
\Large \textbf{Exercise 4: Hyperparameter Tuning} \\[1em]
\normalsize Until lunch
% TODO: add Colab link and QR code
\end{frame}

% ============================================================

\begin{frame}[allowframebreaks]{References}
\printbibliography[heading=none]
\end{frame}

\end{document}
